\subsection{Wait-Awhile under FaaS-Appropriate Parameters}
\label{sec:evaluation:fair-waitawhile}

The headline Table~\ref{tab:headline} uses Wait-Awhile with its
original 1800\,s deferral horizon (Wiesner et al.'s batch-job
default \citep{wiesner2021lets}). A natural reviewer concern is that
1800\,s is adversarial against FaaS functions with sub-minute SLA
deadlines. We address this directly by reporting Wait-Awhile with
shorter deferral horizons matched to FaaS-appropriate timescales.

\begin{table}[h]
\centering\small
\caption{Wait-Awhile under three deferral horizons. Synthetic 24h
workload; real-LWA-carbon results are nearly identical.}
\label{tab:fair-waitawhile}
\begin{tabular}{lrrr}
\toprule
Scheduler                        & Carbon (g) & vs FIFO & Cold rate \\
\midrule
FIFO                             & 51.0      &   0.0\%   & 0.07\% \\
Wait-Awhile (defer=1800\,s)      & 206.6     & $-305.0$\% & 2.02\% \\
Wait-Awhile (defer=60\,s)        & 51.0      &   0.0\%   & 0.07\% \\
Wait-Awhile (defer=30\,s)        & 51.0      &   0.0\%   & 0.07\% \\
\midrule
Spatial                          & 14.0      &  +72.6\%  & 0.03\% \\
\GreenFaaS                       & 14.6      &  +71.4\%  & 0.05\% \\
\bottomrule
\end{tabular}
\end{table}

The empirical finding is more nuanced than the reviewer concern
anticipated. With a FaaS-appropriate deferral horizon
(60\,s or 30\,s), \emph{Wait-Awhile becomes bit-identical to FIFO}:
no future slot meets its absolute-threshold criterion within the
short window, so it always falls back to immediate execution. The
batch-default horizon (1800\,s) produces catastrophic carbon
inflation; the FaaS-appropriate horizons produce no carbon decisions
at all. Wait-Awhile's batch-job design has no parameter sweet spot
for FaaS arrival rates and SLA deadlines.

This finding refines the headline results of \S\ref{sec:evaluation:headline}. The statement ``Wait-Awhile
more than doubles operational carbon on FaaS'' is true under
Wait-Awhile's default parameters; under FaaS-appropriate parameters
the failure mode is different (the scheduler is starved of valid
deferral options and degenerates to FIFO). Either way, Wait-Awhile as
designed does not produce carbon savings on FaaS workloads, but the
mechanism of failure depends on the parameterization: catastrophic at
long horizons, no-op at short horizons. \GreenFaaS's SLA-tiered
deferrable horizon (60\,s for \Deferrable{} class) plus relative
carbon-aware scoring (rather than absolute thresholds) is precisely
what avoids both failure modes.

\paragraph{Threshold parameterization.} The experiments above vary
Wait-Awhile's deferral \emph{horizon} (1800\,s vs 60\,s) while holding
its absolute carbon threshold fixed at $C^\star = 200$\,gCO$_2$/kWh.
A lower threshold would defer fewer invocations (approaching FIFO
sooner) and a higher one would defer more (amplifying the
catastrophic mode at long horizons); neither escapes the underlying
problem, which is that an \emph{absolute} carbon threshold combined
with any fixed horizon cannot adapt to the warm-pool idle-energy
coupling that dominates FaaS. We fix $C^\star$ at a value near the
midpoint of our regions' intensity range so that the comparison
isolates the horizon effect; a full two-dimensional
(horizon $\times$ threshold) sweep is a natural extension but would
not change the qualitative conclusion, since the failure is structural
(absolute thresholding) rather than a matter of threshold tuning.
